\documentclass[11pt]{article}

\usepackage[margin=0.8in]{geometry}
\usepackage{setspace}
\usepackage{parskip}
\usepackage{hyperref}
\usepackage{titlesec}
\usepackage{xcolor}
\usepackage{array}
\usepackage{ragged2e}

\definecolor{PrimaryRed}{HTML}{8B1E1E}
\definecolor{SoftBlue}{HTML}{DDEEFF}
\definecolor{DarkGray}{HTML}{333333}

\hypersetup{
    colorlinks=true,
    linkcolor=black,
    urlcolor=black,
    citecolor=black
}

\titleformat{\section}{\large\bfseries\color{PrimaryRed}}{}{0em}{}
\titleformat{\subsection}{\normalsize\bfseries}{\thesubsection}{0.5em}{}

\setstretch{1.15}

\begin{document}

\begin{titlepage}
\newgeometry{margin=0.85in}
\thispagestyle{empty}

\begin{center}
    \vspace*{0.25in}
    {\large University of Rochester\\}
    {\normalsize Warner School of Education\\[0.35in]}

    {\color{PrimaryRed}\rule{\textwidth}{1.4pt}}\\[0.22in]
    {\Huge\bfseries Joy-Oriented Literacy Pursuit\\[0.12in]}
    {\Large Formal Reflective Journal\\[0.22in]}
    {\color{PrimaryRed}\rule{0.72\textwidth}{0.7pt}}\\[0.32in]

    \colorbox{SoftBlue}{%
        \parbox{0.86\textwidth}{%
            \centering\vspace{0.12in}
            \textbf{Classroom Technology, Design, and Literacy as Joy}\\[0.05in]
            A Teaching Neurodivergent Journal focused on Literature \& Joy via a Puzzle Plan
            \vspace{0.12in}
        }%
    }
\end{center}

\vspace{0.75in}

\begin{center}
\renewcommand{\arraystretch}{1.35}
\begin{tabular}{>{\bfseries\RaggedLeft}p{1.55in} p{4.85in}}
Student & Piter Garcia \\
Course & EDU498: Culturally-Historically Responsive Literacies \\
Class Session & Module 3 / Week 5 \\
Assignment & Joy-Oriented Literacy Pursuit \\
Context & Teaching Reflection through the Lens of this course \& a Puzzle Plan\\
Submission Date & June 05, 2026 \\
\end{tabular}
\end{center}

\vfill

\begin{center}
\begin{minipage}{0.78\textwidth}
\centering
\itshape
Joy-oriented literacy is not extra. It is part of access. When students feel capable, connected, and allowed to participate as themselves, literacy becomes a tool for liberation, self-determination, and empowerment.
\end{minipage}
\end{center}

\restoregeometry
\end{titlepage}
\clearpage

\section*{What I Created and Who I Shared It With}

For my joy-oriented literacy pursuit, I created and facilitated classroom-based technology and designed lessons where students experienced literacy through building, collaboration, speaking, movement, problem-solving, and shared excitement. I focused on creating moments where students could use language, materials, technology, and relationships to make meaning that aligns with the lessons goals. The examples that brought the most joy came from my Kindergarten airplane, light-bulb, and fan circuit lessons, as well as my Minecraft lessons in third (3) and fourth (4) grades. These activities were shared with students in my classroom, but they were also shared with the classroom community itself because students built, watched, helped, responded, revised, and learned from one another in real time.

In my Kindergarten circuit lessons, I gathered students around a demonstration where they learned about batteries, switches, wires, Play-Doh, and motor-based fans by watching, naming parts, holding materials, connecting pieces, and playfully counting down together before turning on a built motor-based fan. These lessons became a shared literacy event because students were participating physically, verbally, visually, and socially, making every concept meaningful. Here, students worked in teams to build an airplane using tools such as Play-Doh and motor-based fans, with different roles such as building wings, shaping the body, making wheels, and connecting the fan to a switch. On the other hand, in my fourth (4) grade Minecraft lessons, students who finished their work were invited to help classmates who were not done, which turned digital literacy into peer teaching and community responsibility. This proves that a shared experience is a literacy product, where students become makers, helpers, designers, and explainers.

\section*{My Creative and Learning Process}

My creative process began with a question: \textit{how can technology lessons become more than just a task completion?} After reading Knobel and Lankshear (2014) and Price-Dennis, Holmes, and Smith (2015), I understood that this question was really about access. Knobel and Lankshear argue that literacy is not fixed to a single form; it evolves as tools, technologies, communities, and purposes change. Price-Dennis, Holmes, and Smith extend this by showing that digital tools only become inclusive when teachers use them intentionally to reduce barriers, increase participation, and support meaningful communication. Together, their work helped me see that digital and media literacy can serve as genuine literacy practices. Students can make meaning through screens, images, videos, comments, games, digital tools, online communities, and physical-digital design. That frame helped me see my lesson activities as literacy work because students were interpreting systems, using vocabulary, communicating with peers, reading visual and technical cues, and representing their thinking through more than one mode.

In the light-bulb and airplane lesson activities, I used materials such as switches, batteries, wires, Play-Doh, fans, and model-building pieces to help students connect vocabulary to action. Students did not only hear the word ``switch''; they saw it, held it, connected it, and watched what happened when it worked. In these lesson activities, students had to decide who would make each part of the airplane, which made the activity collaborative and distributed. Similarly, in one of my fourth (4) grade Minecraft lessons, I used the students who finished early as resources for the class instead of treating completion as the end of learning. That shifted the process from ``I am done'' to ``I can help someone else succeed'', which helps one student to learn and navigate and/or troubleshoot their issues while the other student masters their skills during the process.

This process also taught me to notice joy as part of learning design. Joy did not happen only because the activities were fun. It happened because students had access to materials, roles, language, peer support, and visible success. Knobel and Lankshear make clear that new literacies are social practices built around participation and community, not just technical use of tools. Price-Dennis, Holmes, and Smith reinforce this by arguing that teachers should not simply add a device to a traditional task and call it innovation; instead, the stronger question is whether the task creates more access, voice, connection, and critical thinking. That became the design principle behind this pursuit: \textit{students needed enough structure to feel safe and enough freedom to feel ownership}.

\section*{How This Cultivated Joy, Love, and Aesthetic Fulfillment}

This literacy pursuit cultivated joy because students were able to see their learning become real. In my Kindergarten airplane and fan-circuit lessons, the joy came mostly from collective anticipation. Students watched the pieces connect, counted down together, and saw the fan turn on. That moment mattered because the science was no longer abstract; it became visible, audible, physical, and shared. The joy was not only in the fan working; it was in the class experiencing the success together. These lesson activities created aesthetic fulfillment because students were building something with shape, motion, design, and purpose. Their work now had form, creativity, visible effort, and a relationship to the physical world. The pursuit also cultivated love through relationships. An example is how, in my fourth (4) grade Minecraft lessons, students who already understood the task helped classmates who were still working. That kind of peer support changed the emotional meaning of the activity. Success was not private or competitive; it became something students could give to each other.

These activities reminded me that joy in literacy does not always look quiet or neat. Sometimes joy looks like a student trying again, a group arguing productively over what goes where, a fan finally turning on, a peer stepping in to help, or a student realizing that their hands, ideas, and voice matter. For me, that is aesthetic fulfillment: the moment when learning feels alive, when students can see themselves inside the work, and when the classroom becomes a place where knowledge is built rather than simply delivered.

\section*{Muhammad's Literacy Pursuit and the Module 3 Readings}

\textit{Muhammad's framework} helps me understand that a literacy pursuit should develop identity, skills, intellect, criticality, and joy together. This project was not designed only to teach science skills, but for students to develop vocabulary and technical understanding while also learning how to see themselves as builders, helpers, problem-solvers, and contributors. That matters because literacy should help students understand themselves, participate in community, and experience agency.

\textit{Knobel and Lankshear (2014)} strengthened this connection by helping me name these activities as new literacies, because meaning-making also happens through digital, visual, participatory, and networked forms. Their work challenges the assumption that school-approved print performance is the only valid measure of intelligence or literacy. \textit{Price-Dennis, Holmes, and Smith (2015)} connect the work to inclusive digital literacy by showing that students need more than one pathway to demonstrate understanding. When literacy tasks include multimodal creation, visual design, oral explanation, collaboration, and digital composition, more students can show what they know.

\textit{Brooks and Frankel (2020)} also matter here because students do not enter classrooms as blank technology users. Their work on transnational digital literacy shows that students may already be using digital tools to move across languages, countries, families, and communities before school ever asks them to. When a student builds in Minecraft, explains a circuit orally, draws a design, or helps a peer troubleshoot, that student is using literacy in a way that deserves recognition. These practices are not side activities; they require interpretation, audience awareness, cultural knowledge, and digital navigation.

\textit{Melo-Pfeifer and Gertz (2022)} pushed me to think about criticality. They show that digital and media spaces are not neutral; students need to question headlines, images, missing voices, representation, credibility, and assumptions. The \textit{Common Sense Digital Citizenship} materials extend this by giving students language for privacy, audience, confirmation bias, online identity, community responsibility, and upstander communication. My \textit{Puzzle Plan} project focuses on building the habits that critical digital literacy needs: \textit{collaboration, questioning, explanation, audience awareness, and responsibility to others}.

\textit{Stornaiuolo and Thomas (2017)} helped me understand agency and joy. They show that young people can use digital tools to speak, organize, and participate in public life---positioning students not as passive technology users but as active participants who create public meaning. \textit{Pinkard's ecosystem} view of digital literacy reminds me that students need tools, mentors, audiences, feedback, and real reasons to create. The Grade 4 Minecraft peer-teaching moment was exactly that ecosystem in action. \textit{Serpagli and Mensah's (2021)} work with social media-style science learning connects directly to my placement examples because students were using familiar, active, and shared forms to explain technical ideas. Together, these readings show that digital literacy should position students as creators, not just consumers.

\section*{How This Informs My Future Teaching}

This experience informs my teaching by reminding me that literacy is not limited to written language. Literacy is also building, coding, designing, explaining, translating, collaborating, troubleshooting, helping, evaluating media, and trying again. In my classroom, I will design activities where students have multiple entry points and different ways to show their understanding. This is especially important for students who may be multilingual, neurodivergent, disabled, chronically ill, culturally diverse, or simply different from the ``default'' learner that schools often imagine.

In practice, I will continue designing projects where students create something meaningful and share it with an authentic audience. Students could build Minecraft worlds connected to a science concept, create bilingual or multimodal explanations, make collaborative classroom murals, create technology tutorials for peers, or build simple circuitry projects and explain how they work through speaking, drawing, writing, or video. I will also make peer teaching part of the classroom culture because students often understand each other's barriers in ways adults often miss.

My biggest takeaway is that joy-oriented literacy is part of access. Through my \textit{Puzzle Plan and EQUITAS lens}, the question is not whether a task looks traditional enough to count as literacy. The question is whether the task gives students real access to meaning, identity, criticality, collaboration, and agency. Knobel and Lankshear remind us that literacy changes as communities and tools change; Price-Dennis and colleagues show that inclusion requires intentional design; Brooks and Frankel show that students' existing digital and cultural practices are already literacy; and Stornaiuolo and Thomas show that digital tools can be instruments of voice, not just task completion. When students feel capable, connected, and allowed to participate as themselves, literacy becomes a tool for liberation, self-determination, and empowerment. This pursuit helped me see that my teaching practice should ask whether students had the chance to feel ownership, joy, relationship, and possibility.

\end{document}
